Despite advances in drought monitoring, significant limitations remain in translating meteorological conditions into actionable impact predictions. Current monitoring systems in the Netherlands are primarily based on physical indicators such as precipitation-based indices (e.g., Standardized Perception index) \cite{KNMI2023}. While these indices effectively capture anomalies in climate variables, they provide limited insight into the sector-specific consequences of drought events \cite{Shyrokaya2024Advances}.

Drought impacts extend across multiple domains. In the Netherlands, long dry periods can lead to groundwater depletion, reduced agricultural yields, ecological stress and increased salinisation in coastal and delta regions \cite{KNMI2023}. These impacts illustrate that drought risk is not solely a meteorological phenomenon but a complex socio-hydrological process.

Recent literature on drought impact modelling has increasingly focused on predictive approaches using existing datasets such as the European Drought Impact Inventory (EDII). However, a key limitation of these datasets is that they were not originally designed for impact-based forecasting.  As a result, they often have limited spatial resolution, inconsistent reporting standards \cite{Shyrokaya2024Advances}. While countries such as Germany and the United Kingdom, and Austria have developed methodologies for drought impact analysis and prediction, these efforts remain constrained by data quality and granularity, which limits their applicability for high-resolution forecasting tasks \cite{Bulut2026Framework, Sutanto2019Moving}.

In the Netherlands, there is currently no comprehensive national-scale drought impact database or operational impact-based forecasting framework. The Dutch section of the EDID2 database contains only 219 reported drought impacts, compared with 1,855 reports for the United Kingdom. In contrast, the forecasting framework developed by Bulut was trained using 4,738 drought impact reports  \cite{Bulut2026Framework}. 

% This gap is particularly notable given the country's strong dependence on engineered water management systems, as large parts of the territory require active water regulation because they lie below sea level. Consequently, the Netherlands represents a particularly relevant case for the development of fine-grained, data-driven approaches to drought impact modelling.


This thesis identifies two gaps that limit the development of impact-based drought forecasting systems:

\begin{itemize}

\item \textbf{(1) Drought-impact database gap}: There is currently no structured, multi-sectoral drought impact database at the national scale for the Netherlands that enables systematic analysis of drought-related impacts.

\item \textbf{(2) Forecasting framework gap:} The availability of a national drought impact database creates new opportunities for impact-based forecasting. However, no baseline forecasting framework currently exists for the Netherlands to assess the predictability of drought impacts or to establish the inherent performance limits of such models.

\end{itemize}


To address these gaps, this thesis investigates how a national drought impact database can be constructed and used for impact-based drought forecasting.

The thesis consists of two parts. Part I develops a geospatial drought impact database from newspaper articles using large language models (LLMs) and evaluates the quality of the extracted impact records. Part II uses this database to develop and evaluate impact-based drought forecasting models based on a theoretical framework integrating hazard, exposure, and impact-derived features. Together, these components provide a foundation for future operational impact-based forecasting.
